% Source: source/普通高考/2013/2013四川文.pdf, page 1 (lower middle-right option D).
% The source PDF has no embedded raster objects (pdfimages -list is empty); the
% chart is vector linework.  A 100-unit grid inspection of a 300 dpi render
% identified four adjacent bins of width 10 on [0,40], with heights
% 0.015, 0.03, 0.035, 0.02.  Bars share edges and rise from the baseline;
% axes have arrowheads, y ticks are 0.01--0.04, and labels are “频率/组距” and
% “人数”.
\begin{tikzpicture}
  \begin{axis}[exam statistical histogram,
    width=5.8cm,height=4.0cm,
    xmin=0,xmax=42,ymin=0,ymax=0.052,
    axis lines=left, axis line style={->},
    xtick={0,10,20,30,40},
    xticklabel=\empty,
    ytick={0.01,0.02,0.03,0.04}, scaled y ticks=false,
    yticklabel style={font=\normalsize,/pgf/number format/fixed,/pgf/number format/precision=2},
    tick style={black}, tick label style={font=\normalsize},
    ylabel={},
    tick align=outside, clip=false,
  ]
    \examstatxlabel{人数};
    \examstatylabel{\(\dfrac{\text{频率}}{\text{组距}}\)};
    \node[anchor=north east,inner sep=0pt,xshift=-0.5pt,font=\normalsize] at (rel axis cs:0,-0.035) {0};
    \node[anchor=north,inner sep=0pt,font=\normalsize] at (rel axis cs:0.238095,-0.035) {10};
    \node[anchor=north,inner sep=0pt,font=\normalsize] at (rel axis cs:0.476190,-0.035) {20};
    \node[anchor=north,inner sep=0pt,font=\normalsize] at (rel axis cs:0.714286,-0.035) {30};
    \node[anchor=north,inner sep=0pt,font=\normalsize] at (rel axis cs:0.952381,-0.035) {40};
    \addplot[/tikz/ybar interval,draw=black,fill=white,mark=none] coordinates {
      (0,0.015) (10,0.03) (20,0.035) (30,0.02) (40,0)
    };
  \end{axis}
\end{tikzpicture}
